% Recommended AQT main-text figure set.

\begin{figure*}[t]
\centering
\includegraphics[width=.98\textwidth]{fig1_framework_theory.pdf}
\caption{Measurement-accessible parameter-tangent geometry. A fixed quantum measurement defines a centered score space. The tangent-score covariance $C$ distributes normalized parameter-tangent mass within that space, while a rank-$r$ readout projector $P$ selects the retained directions. The mean retained mass is $R=\operatorname{Tr}(PC)$. Random relative orientation supplies the rank-only reference $r/N$; the covariance spectrum sets the corresponding null width, and the physical orientation of $P$ determines the observed deviation.}
\label{fig:framework-aqt}
\end{figure*}

\begin{figure*}[t]
\centering
\includegraphics[width=.98\textwidth]{fig2_rank_architecture_anisotropy.pdf}
\caption{Rank-normalized retention and covariance anisotropy. Family-balanced generic low-weight readouts remain close to the random-orientation rank reference over the tested sizes while the normalized purity $N\operatorname{Tr}(C^2)$ grows and $d_{\mathrm{eff}}/N$ decreases. Architecture-resolved deviations show that near-rank aggregate retention is not evidence of isotropic or Haar-oriented tangent geometry.}
\label{fig:rank-aqt}
\end{figure*}

\begin{figure*}[t]
\centering
\includegraphics[width=.98\textwidth]{fig4_spectral_mechanism_depth.pdf}
\caption{Equal-rank spectral comparison. Physical low-weight, cross-fitted aligned, and same-sample Ky Fan subspaces have the same retained dimension but can capture different amounts of tangent mass. The orientation gap is large in generic Haar-$U(4)$ circuits and smaller in the half-filled $U(1)$ family, where the physical low-weight readout is already closer to the leading tangent subspace.}
\label{fig:spectrum-aqt}
\end{figure*}

\begin{figure*}[t]
\centering
\includegraphics[width=.95\textwidth]{appendix_operational_gains.pdf}
\caption{Finite-shot consequence of equal-rank readout orientation. The panels compare the physical and cross-fitted aligned readouts at fixed circuit, computational-basis measurement, readout rank, evaluation directions, and shot budget. Generic Haar-$U(4)$ circuits show increasing directional gradient-energy and signal-to-noise gains with system size, whereas the structured $U(1)$ readout leaves less room for orientation improvement.}
\label{fig:operational-aqt}
\end{figure*}

\begin{figure*}[t]
\centering
\includegraphics[width=.95\textwidth]{production/fig_u1_sector_null.pdf}
\caption{Sector-corrected random-orientation control for the half-filled $U(1)$ alignment statistic. The null draws the leading rank-$r_1$ subspace randomly inside the same fixed-charge centered score sector. Measured cross-fitted low-weight alignment remains far above this sector-corrected reference at the largest sizes, showing that the structured result is not produced by support or rank reduction alone.}
\label{fig:u1-aqt}
\end{figure*}
